\section{Introduction}

Cell-free massive multiple-input multiple-output (MIMO) has emerged as a promising architecture for 6G wireless networks. By distributing a large number of access points (APs) over a wide geographical area to coherently serve multiple user equipments (UEs), cell-free MIMO effectively eliminates cell boundaries, thereby providing substantial macro-diversity and spatial multiplexing gains \cite{bashar2020uplink, park2025scalable}. In this architecture, the distributed APs are connected to a central processing unit (CPU) via fronthaul or backhaul links, enabling joint signal processing and interference management \cite{li2024joint, park2013multihop}. However, the performance of cell-free MIMO systems heavily relies on the capacity of these fronthaul and backhaul links. As the network scales in terms of the number of APs and antennas, the massive volume of data exchanged between the APs and the CPU imposes a fundamental performance bottleneck, especially when the fronthaul capacity is strictly limited and dynamically fluctuating \cite{bashar2020exploiting, ranjbar2024joint}.

In conventional centralized processing schemes, APs forward their received signals to the CPU using high-resolution quantization (e.g., 32-bit floating-point representation). While this preserves signal integrity, it requires significant fronthaul bandwidth, which can be prohibitive for dense deployments \cite{park2025scalable}. To alleviate this burden, various low-bit quantization and compression techniques have been proposed. Traditional approaches often rely on uniform fixed bit-width quantization or heuristic bit-allocation algorithms \cite{bashar2020uplink, kim2024meta}. However, these methods present certain limitations. First, they typically apply a static quantization resolution across all APs, failing to adapt to instantaneous channel state information (CSI) and dynamic fronthaul bandwidth availability. Second, conventional quantization schemes are generally decoupled from the end-to-end receiver task (e.g., signal detection). Consequently, they cannot perform task-oriented optimal bit allocation, meaning they are unable to prioritize the transmission of the most critical demodulation information under severe bandwidth constraints, leading to potential resource inefficiencies or degradation in detection accuracy \cite{ranjbar2024joint}.

Recently, artificial intelligence (AI) and deep learning have been applied to physical layer communications to address the limitations of traditional algorithms. Model-driven deep learning has demonstrated efficacy in MIMO detection by unfolding iterative algorithms to reduce computational complexity while improving robustness \cite{he2020model, wei2020learned}. Furthermore, graph neural networks (GNNs) have been increasingly adopted in cell-free MIMO systems for decentralized processing, resource allocation, and interference coordination, owing to their scalability and permutation invariance \cite{pereira2025pixel, wang2025learning, dai2025two}. Unsupervised and meta-learning approaches have also been explored to optimize power control and signal compression under fronthaul constraints \cite{rajapaksha2023unsupervised, xie2026meta}. Despite these advancements, a critical research gap remains: existing AI-based models rarely address the non-differentiable nature of dynamic, discrete bit-width allocation under strict bandwidth constraints. Most learning-based quantization schemes either assume continuous capacity relaxations or fail to jointly optimize the discrete quantization policy with the non-linear spatial interference cancellation at the CPU. Consequently, how to design a fully differentiable, task-oriented successive refinement framework that autonomously infers the optimal discrete bit-width for each AP based on local channel conditions remains an open challenge.

To bridge this gap, this paper proposes a resource-inference GNN framework for backhaul-constrained cell-free MIMO systems based on a task-oriented successive refinement approach. By integrating a progressive quantizer with a learnable resource-gating module, the proposed system dynamically allocates discrete bit-widths to APs, ensuring that the most critical information is preserved even under severe bandwidth limitations. The main contributions of this paper are summarized as follows:
\begin{itemize}
    \item We design a differentiable Learned Step Size Quantization (LSQ) module that enables successive refinement of the received signals at the APs. This allows the continuous baseband signals to be decomposed into hierarchically ordered bit layers, facilitating flexible truncation based on instantaneous bandwidth constraints.
    \item We propose a lightweight, Gumbel-Softmax-driven bit-width policy network deployed at the APs for autonomous resource inference. By transforming the discrete bit-width selection into a differentiable operation during training, the policy network learns to dynamically allocate bit-widths (e.g., 0, 2, and 4 bits) based solely on local channel conditions, effectively discarding severely faded signals to conserve backhaul resources.
    \item We develop a joint Quantization-Aware Training (QAT) framework that integrates the distributed policy networks with a heterogeneous GNN at the CPU. The GNN employs a bit-aware attention mechanism to adaptively process signals of asymmetric precision, performing robust spatial interference cancellation and non-linear detection.
    \item We conduct extensive numerical evaluations demonstrating the robustness and scalability of the proposed framework. Under heavy user load scenarios, the proposed method reduces the bit error rate (BER) by 67.3\% compared to conventional fixed 2-bit quantization schemes. Furthermore, the system exhibits resilience to link failures, maintaining a BER of less than 2.2\% even when 50\% of the APs experience complete backhaul dropout. Additionally, in sparse coverage scenarios with severe path-loss, the policy network autonomously compresses the average bit-width to discard heavily faded signals. Finally, complexity analysis reveals that the AP-side policy network maintains a constant computational overhead (e.g., 384 floating-point operations) entirely independent of the antenna array size, ensuring scalability for massive MIMO deployments.
\end{itemize}

The remainder of this paper is organized as follows. Section II systematically defines the system model, including the physical communication environment, the limited-fronthaul cell-free MIMO architecture, and the mathematical problem formulation. Section III details the proposed task-oriented successive refinement methodology, encompassing the differentiable quantization module, the Gumbel-Softmax-driven policy network, and the joint QAT GNN framework. Section IV presents extensive numerical results and physical insights to validate the performance, robustness, and scalability of the proposed approach. Finally, Section V concludes the paper and discusses potential directions for future research.